% ════════════════════════════════════════════════════════════════
%  术语约定与对照（frontmatter，全书统一术语口径）
%  注：本书集机器人学多子领域之大成，部分术语在不同方向有不同（且各自正确）的叫法。
%      本表分两类登记：(1)「按方向区分」——同形不同义 / 同义不同名，须按语境择用，切勿一刀切统一；
%      (2)「全书统一」——单一首选译名（同义词作备注）。本表是各章用词的约定基准，亦防多源融合时的术语漂移。
% ════════════════════════════════════════════════════════════════
\chapter*{术语约定与对照}
\addcontentsline{toc}{chapter}{术语约定与对照}
\label{ch:terminology}
\markboth{术语约定与对照}{术语约定与对照}

本书融汇估计、SLAM、规划、控制、强化学习、足式与操作等多个子领域。\textbf{同一英文术语在不同方向常有不同且各自正确的中文叫法，亦有同形中文词在不同方向指代不同概念}。强行抹平这些差异反而会引入错误。故本表分两类登记：

\begin{itemize}
  \item \textbf{按方向区分}（见下\,\textbf{表一}）：同形不同义、或同义不同名，\emph{须按语境择用}。例如“回报”（return $G_t$）与“奖励”（reward $R_t$）是强化学习中\emph{两个不同概念}，不可混同；“构型/位形”（关节/广义坐标）与“位姿”（$\SE(3)$ 刚体）各有所指；“镇定”（stabilization）与“稳定”（stability）近义而不同。
  \item \textbf{全书统一}（见下\,\textbf{表二}）：同一概念的多个中文译名中，本书择一为首选，其余作同义词备注，供检索与对照。
\end{itemize}

各章末“符号表”与本书《主要符号与记号约定》（\cref{ch:notation}）登记\emph{符号}；本表登记\emph{术语}，二者互补。

\section*{一、按方向区分（须按语境择用，勿一刀切）}

\begin{center}\small
\begin{longtable}{p{0.20\textwidth} p{0.34\textwidth} p{0.36\textwidth}}
\caption*{\textbf{表一}\quad 同形不同义 / 同义不同名的术语：按子领域语境区分}\\
\toprule
英文 / 概念 & 本书用词（按方向） & 区分要点 \\
\midrule
\endfirsthead
\multicolumn{3}{c}{\small（续）术语按方向区分}\\
\toprule
英文 / 概念 & 本书用词（按方向） & 区分要点 \\
\midrule
\endhead
\bottomrule
\endfoot
pose vs.\ configuration
  & \textbf{位姿}（pose，刚体）／\textbf{构型}（configuration，关节空间）
  & 位姿 = 刚体在 $\SE(3)$ 的位置+姿态（6 DOF）；构型 = 广义坐标 $\mathbf{q}$（关节角等）所定的机构组态。两者\emph{不同概念}，勿混。\textbf{构型 vs.\ 位形}：“构型（空间）”是 GB/T 12643 国标与 \emph{Modern Robotics} 用词、\textbf{本书首选}；“位形（空间）”是 Barfoot 等物理传统的等价说法，本书亦见，二者所指相同。\\
reward vs.\ return
  & \textbf{奖励}（$R_t$）／\textbf{回报}（$G_t$）
  & 奖励 = MDP 每步的即时反馈标量 $R_t$；回报 = 从某时刻起的（折扣）累积奖励 $G_t=\sum_k\gamma^k R_{t+k}$。强化学习中是\emph{两个不同量}，不可互换。\\
controllability / observability
  & 控制论：\textbf{能控性 / 能观性}；估计·SLAM：\textbf{可观性}
  & 经典控制（状态空间综合，承刘豹《现代控制理论》脉络，见控制部）沿用“能控/能观”；状态估计与 SLAM 的“可观性分析”（哪些状态可由观测恢复）用“可观性”。同一英文、不同传统，按所在子领域用词。\\
closed-loop vs.\ loop closure
  & 控制：\textbf{闭环}；SLAM：\textbf{回环}
  & 闭环 = 反馈控制回路（closed-loop）；回环 = SLAM 重访旧地形成的约束（loop closure）。中文同含“环”，含义\emph{无关}，勿互换。\\
旋量族 (screw / twist / wrench)
  & \textbf{螺旋轴}（screw axis $\mathcal{S}$）／\textbf{运动旋量}（twist $\mathcal{V}$）／\textbf{力旋量}（wrench $\mathcal{F}$）
  & “旋量”是旋量理论的\emph{根名词}，单独使用时含义\textbf{随语境}：POE/旋量理论中指（单位）\textbf{螺旋轴} $\mathcal{S}=[\boldsymbol{\omega};\mathbf{v}]$（本书亦记“螺旋轴”）；运动学中的速度六元（角速度+线速度）为\textbf{运动旋量} $\mathcal{V}$（twist，亦称速度旋量）；其对偶（力+力矩）为\textbf{力旋量} $\mathcal{F}$（wrench）。\textbf{跨源警示}：\emph{Modern Robotics} 官方中译本中“旋量”\emph{单独即指 twist}，而本书 POE 语境的“旋量”指 screw axis——援引不同来源务必按上下文判别，勿把两义混用。排序见 \cref{ch:notation}。\\
stabilization vs.\ stability
  & \textbf{镇定}（stabilization）／\textbf{稳定性}（stability）
  & 镇定 = 控制\emph{目标/动作}（设计反馈使平衡点渐近稳定，“要做什么”）；稳定性 = 系统\emph{性质}（“系统表现如何”）。近义但\emph{不同}：可镇定 $\ne$ 已稳定。\\
pitch（三义）
  & 语境决定：\textbf{螺距}（screw pitch $h$）／\textbf{俯仰角}（pitch angle）／螺纹螺距
  & 同一英文“pitch”三个独立概念：旋量理论的螺距 $h=$ 线速度/角速度（标量，量纲长度）；姿态欧拉角的俯仰角（角度）；螺纹的螺距（长度）。本书主用俯仰角；涉旋量理论处须显式标“螺距”以免混。\\
Hessian vs.\ Hamiltonian
  & \textbf{海森}（Hessian 矩阵）／\textbf{哈密顿量}（Hamiltonian）
  & 英文形近、中文音近但\emph{毫不相干}：海森 $=$ 二阶偏导矩阵 $\partial^2 f/\partial\mathbf{x}^2$；哈密顿量 $=$ 最优控制/力学的能量型函数 $H$（PMP、Poisson 括号）。\\
spatial frame
  & \textbf{基座系 / 世界系}（固定参考），\emph{非}“空间系”
  & “spatial”英文指代模糊（物理空间 vs 固定于空间的参考系）；本书\emph{避用}“空间系”（易误读为“空间维度/向量空间”），固定参考系一律用“世界系”（惯性）或“基座系”（浮动基），随动系用“物体系/机身系”。\\
\end{longtable}
\end{center}

\section*{二、全书统一（首选译名，括注同义词）}

\begin{center}\small
\begin{longtable}{p{0.26\textwidth} p{0.20\textwidth} p{0.44\textwidth}}
\caption*{\textbf{表二}\quad 全书统一首选译名与同义词}\\
\toprule
英文 & 本书首选 & 同义词 / 备注 \\
\midrule
\endfirsthead
\multicolumn{3}{c}{\small（续）全书统一译名}\\
\toprule
英文 & 本书首选 & 同义词 / 备注 \\
\midrule
\endhead
\bottomrule
\endfoot
bundle adjustment & 束调整 & $\approx$ 光束法平差（传统测量学术语）；缩写 BA \\
value function & 值函数 & 取 Sutton--Barto 中译惯例；避免与日常“价值”混淆 \\
observability（估计/SLAM） & 可观性 & $\approx$ 可观测性；与控制部“能观性”为同概念、不同子领域用词 \\
Hessian & Hessian & 亦作“海森（矩阵）”；与 Hamilton 量（哈密顿量）区分 \\
Jacobian & 雅可比 & 雅可比矩阵 \\
manifold & 流形 & $\SO(3)/\SE(3)$ 等微分几何对象 \\
residual & 残差 & 观测与模型预测之差 \\
covariance & 协方差 & 矩阵形式为协方差矩阵 \\
marginalization & 边缘化 & 因子图/概率图中消元，对应 Schur 补 \\
impedance / admittance & 阻抗 / 导纳 & 对偶概念（力$\to$位移 / 位移$\to$力）\\
compliance & 柔顺 & 刚度之逆，侧重静态顺应 \\
underactuated & 欠驱动 & 独立控制输入数 $<$ 自由度数 \\
nonholonomic & 非完整 & 速度级不可积约束（Pfaffian）\\
attitude / orientation & 姿态 & 刚体朝向；“朝向”为同义词 \\
descriptor & 描述子 & 特征描述子 \\
keyframe & 关键帧 & \\
odometry & 里程计 & 视觉/激光/惯性里程计 \\
policy & 策略 & 强化学习中的 $\pi(a\mid s)$ \\
configuration & 构型 & GB/T 12643 国标用词；$\approx$ 位形（Barfoot 物理传统）。“构型空间/位形空间”即 $\mathcal{C}$-space \\
innovation（Kalman） & 新息 & 实测与预测之差（新的信息）；\emph{勿}译“创新”（那是 novelty 本义）\\
rotation matrix & 旋转矩阵 & 本书统一“旋转矩阵”，不用“转动矩阵”；亦称方向余弦矩阵 \\
manipulability / dexterity & 可操作度 / 灵巧度 & 可操作度 $=$ Yoshikawa $w=\sqrt{\det(\mathbf{J}\mathbf{J}^\top)}$（运行时）；灵巧度 $=$ 设计级指标（条件数/GCI）\\
redundancy & 冗余 & 自由度 $>$ 任务维时的结构性质；\emph{不}用“冗余度” \\
regulation / tracking & 调节 / 跟踪 & 调节 $=$ 镇定到固定目标；跟踪 $=$ 跟随时变参考（反馈+前馈）\\
force / form closure & 力封闭 / 形封闭 & 力封闭含摩擦、抗任意力旋量；形封闭纯几何约束（无摩擦）\\
end-effector & 末端执行器 & ISO 标准术语；亦简称“末端” \\
gimbal lock & 万向锁 & 欧拉角奇异（俯仰 $=\pm\pi/2$ 丢一自由度）；勿与硬件万向节故障混 \\
\end{longtable}
\end{center}

\begin{center}\small
\emph{说明}：本表条目经\textbf{双盲联网核对}（多源独立核验 $+$ 裁决，锚 \emph{Modern Robotics}/Craig/Spong/Siciliano/Sutton--Barto 官方中译本、GB/T 12643《机器人术语》国标、维基中英对照）。本表随全书演进增补：新增章节用词应先查本表；若遇本表未列、且不同方向叫法有别的术语，按“同子领域既有多数派”择用，并回补本表。
\end{center}
